% =====================================================================
%  LAMPIRAN A — CONTOH LAMPIRAN
%  Bab ini hanya contoh singkat. \appendix pada main.tex mengubah
%  penomoran bab menjadi huruf (LAMPIRAN A) dan penomoran objek menjadi
%  A.1, A.2, dan seterusnya.
% =====================================================================

\chapter{Data Pendukung}
\label{lam:a}

\section{Contoh Tabel dan Persamaan}
\label{lam:a.1}

[Lampiran ini menunjukkan cara menempatkan tabel, persamaan, dan rujukan
silang di dalam lampiran. Ganti isinya dengan data pendukung penelitian
Anda, misalnya tabel pengukuran lengkap, kuesioner, atau kode program.]

Tabel~\ref{tab:lam.a1} merangkum contoh data pendukung.

\begin{thesistable}
\begin{longtable}{@{}l r r@{}}
  \caption{Contoh data pendukung pada lampiran}
  \label{tab:lam.a1}\\
  \toprule
  \textbf{Pengamatan} & \textbf{[Nilai 1]} & \textbf{[Nilai 2]} \\
  \midrule
  \endfirsthead
  \multicolumn{3}{@{}l}{\footnotesize\itshape Lanjutan Tabel~\thetable}\\
  \toprule
  \textbf{Pengamatan} & \textbf{[Nilai 1]} & \textbf{[Nilai 2]} \\
  \midrule
  \endhead
  \bottomrule
  \endfoot
  \bottomrule
  \endlastfoot
  {[Pengamatan 1]} & [X,XX] & [X,XX] \\
  {[Pengamatan 2]} & [X,XX] & [X,XX] \\
  {[Pengamatan 3]} & [X,XX] & [X,XX] \\
\end{longtable}
\end{thesistable}

Rata-rata dan varians contoh data ditunjukkan pada Persamaan~\eqref{eq:lam.a1}
dan Persamaan~\eqref{eq:lam.a2}.

\begin{equation}
  \bar{x} = \frac{1}{n}\sum_{i=1}^{n} x_i
  \label{eq:lam.a1}
\end{equation}

\begin{equation}
  s^2 = \frac{1}{n-1}\sum_{i=1}^{n}\left(x_i - \bar{x}\right)^2
  \label{eq:lam.a2}
\end{equation}
